% GENERATED FILE. Edit canonical lessons, then run npm run generate:books.
% canonical-source-hash: fnv1a64:8af58e87c1215032
% canonical-lessons: GE-C501-brettspiel, GE-C501-schach, GE-C501-wuerfel, GE-C501-mogeln, GE-C501-runde

\chapter{Brettspiele}
\label{ch:ge-brettspiele}
\hlchaptermodality{501}

\begin{chapteropening}
\textbf{By the end of this chapter:} \emph{I can talk about board games.}

Complete the last lesson of chapter 501: I can talk about board games.
\end{chapteropening}

\section[das Brettspiel]{das Brettspiel — board game}

\label{lesson:GE-C501-brettspiel}



\begin{quote}
Before the new one: say the German for riddle, then the German for crossword.
\end{quote}

\subsection*{You'll want to know: das Brettspiel}
\textbf{das Brettspiel} — \textquotedblleft{}board game\textquotedblright{}.

\begin{quote}
Am Wochenende spielen wir gern Brettspiele.
\end{quote}

\begin{quote}
\emph{At the weekend we like playing board games.}
\end{quote}

Plural: \emph{die Brettspiele}.

\begin{cousinweb}
\emph{Brettspiel} is built on \emph{Brett} \textquotedblleft{}board, plank\textquotedblright{}, inherited from Middle High German \emph{bret}, and \emph{Spiel} \textquotedblleft{}game\textquotedblright{}, from Old High German \emph{spil} (EtymDB).
\end{cousinweb}

\begin{grammarlens}[title={What you've built}]
One more word for telling, explaining and describing at B1.
\end{grammarlens}

\subsection*{Your turn}
\begin{itemize}
\raggedright
  \item \emph{Say it:} \emph{das Brettspiel}
  \item \emph{Say it:} \emph{das Brettspiel}, once more
  \item \emph{Say it:} a sentence of your own with \emph{das Brettspiel}
  \item \emph{Recall:} say the German for riddle, then the German for crossword, then say \emph{das Brettspiel} again
\end{itemize}

\begin{tcolorbox}[breakable,colback=teal!4,colframe=teal!35!black,title={Before you move on}]
What does das Brettspiel mean? (\textquotedblleft{}Board game\textquotedblright{}.) Say it once more.
\end{tcolorbox}
\section[das Schach]{das Schach — chess}

\label{lesson:GE-C501-schach}



\begin{quote}
Before the new one: say the German for crossword, then the German for board game.
\end{quote}

\subsection*{You'll want to know: das Schach}
\textbf{das Schach} — \textquotedblleft{}chess\textquotedblright{}.

\begin{quote}
Mein Sohn spielt jeden Dienstag Schach im Verein.
\end{quote}

\begin{quote}
\emph{My son plays chess at the club every Tuesday.}
\end{quote}

No plural. \emph{Schach!} — check! \emph{Schachmatt} — checkmate.

\begin{cousinweb}
\emph{Schach} came through Medieval Latin \emph{scaccus} from Arabic, which took it from Persian \emph{shāh} \textquotedblleft{}king\textquotedblright{} (EtymDB): the game is named after its key piece.
\end{cousinweb}

\begin{grammarlens}[title={What you've built}]
One more word for telling, explaining and describing at B1.
\end{grammarlens}

\subsection*{Your turn}
\begin{itemize}
\raggedright
  \item \emph{Say it:} \emph{das Schach}
  \item \emph{Say it:} \emph{das Schach}, once more
  \item \emph{Say it:} a sentence of your own with \emph{das Schach}
  \item \emph{Recall:} say the German for crossword, then the German for board game, then say \emph{das Schach} again
\end{itemize}

\begin{tcolorbox}[breakable,colback=teal!4,colframe=teal!35!black,title={Before you move on}]
What does das Schach mean? (\textquotedblleft{}Chess\textquotedblright{}.) Say it once more.
\end{tcolorbox}
\section[der Würfel]{der Würfel — dice; cube}

\label{lesson:GE-C501-wuerfel}



\begin{quote}
Before the new one: say the German for board game, then the German for chess.
\end{quote}

\subsection*{You'll want to know: der Würfel}
\textbf{der Würfel} — \textquotedblleft{}dice; cube\textquotedblright{}.

\begin{quote}
Wer mit dem Würfel eine Sechs wirft, darf anfangen.
\end{quote}

\begin{quote}
\emph{Whoever throws a six with the dice gets to start.}
\end{quote}

Plural: \emph{die Würfel}. Also a cube: \emph{der Eiswürfel} — an ice cube. The verb is \emph{würfeln} — to throw dice.

\begin{cousinweb}
\emph{Würfel} is inherited from Middle High German \emph{würfel} and Old High German \emph{wurfil}, from Proto-Germanic \emph{wurpilaz} \textquotedblleft{}die\textquotedblright{} (EtymDB), built on the verb that is today \emph{werfen} \textquotedblleft{}to throw\textquotedblright{}.
\end{cousinweb}

\begin{grammarlens}[title={What you've built}]
One more word for telling, explaining and describing at B1.
\end{grammarlens}

\subsection*{Your turn}
\begin{itemize}
\raggedright
  \item \emph{Say it:} \emph{der Würfel}
  \item \emph{Say it:} \emph{der Würfel}, once more
  \item \emph{Say it:} a sentence of your own with \emph{der Würfel}
  \item \emph{Recall:} say the German for board game, then the German for chess, then say \emph{der Würfel} again
\end{itemize}

\begin{tcolorbox}[breakable,colback=teal!4,colframe=teal!35!black,title={Before you move on}]
What does der Würfel mean? (\textquotedblleft{}Dice; cube\textquotedblright{}.) Say it once more.
\end{tcolorbox}
\section[mogeln]{mogeln — to cheat (in a game)}

\label{lesson:GE-C501-mogeln}



\begin{quote}
Before the new one: say the German for chess, then the German for dice.
\end{quote}

\subsection*{You'll want to know: mogeln}
\textbf{mogeln} — \textquotedblleft{}to cheat (in a game)\textquotedblright{}.

\begin{quote}
Beim Kartenspielen mogelt mein kleiner Bruder immer.
\end{quote}

\begin{quote}
\emph{My little brother always cheats at cards.}
\end{quote}

Informal; also \emph{schummeln}.

\begin{cousinweb}
The origin of \emph{mogeln} is not settled; EtymDB lists the German word without a source.
\end{cousinweb}

\begin{grammarlens}[title={What you've built}]
One more word for telling, explaining and describing at B1.
\end{grammarlens}

\subsection*{Your turn}
\begin{itemize}
\raggedright
  \item \emph{Say it:} \emph{mogeln}
  \item \emph{Say it:} \emph{mogeln}, once more
  \item \emph{Say it:} a sentence of your own with \emph{mogeln}
  \item \emph{Recall:} say the German for chess, then the German for dice, then say \emph{mogeln} again
\end{itemize}

\begin{tcolorbox}[breakable,colback=teal!4,colframe=teal!35!black,title={Before you move on}]
What does mogeln mean? (\textquotedblleft{}To cheat (in a game)\textquotedblright{}.) Say it once more.
\end{tcolorbox}
\section[die Runde]{die Runde — round}

\label{lesson:GE-C501-runde}



\begin{quote}
Before the new one: say the German for dice, then the German for to cheat.
\end{quote}

\subsection*{You'll want to know: die Runde}
\textbf{die Runde} — \textquotedblleft{}round\textquotedblright{}.

\begin{quote}
Nach drei Runden hatte meine Schwester gewonnen.
\end{quote}

\begin{quote}
\emph{After three rounds my sister had won.}
\end{quote}

Plural: \emph{die Runden}. \emph{Eine Runde ausgeben} — to buy a round of drinks.

\begin{cousinweb}
\emph{Runde} is built on \emph{rund} \textquotedblleft{}round\textquotedblright{}, which came through Middle High German \emph{runt} from Old French \emph{ront} (EtymDB), itself from Latin \emph{rotundus}.
\end{cousinweb}

\begin{grammarlens}[title={What you've built}]
One more word for telling, explaining and describing at B1.
\end{grammarlens}

\subsection*{Your turn}
\begin{itemize}
\raggedright
  \item \emph{Say it:} \emph{die Runde}
  \item \emph{Say it:} \emph{die Runde}, once more
  \item \emph{Say it:} a sentence of your own with \emph{die Runde}
  \item \emph{Recall:} say the German for dice, then the German for to cheat, then say \emph{die Runde} again
\end{itemize}

\begin{tcolorbox}[breakable,colback=teal!4,colframe=teal!35!black,title={Before you move on}]
What does die Runde mean? (\textquotedblleft{}Round\textquotedblright{}.) Say it once more.
\end{tcolorbox}
